% Template of Appendix_Gemma.tex (Online Appendix). Filled in by 05_appendix_gemma.ipynb;
% <<name>> placeholders are replaced by values computed from the estimation outputs.
\hyphenation{EmbeddingGemma}
\begingroup
\setlength{\emergencystretch}{3em}% the model name is not hyphenated
\section{Robustness Check: Embeddings without Fine-Tuning}\label{app:gemma}

\revised{This section repeats the analysis with off-the-shelf embeddings that are not fine-tuned. The workflow is exactly the same as in the main analysis (Section~\ref{app:workflow}), except that the fine-tuning in step~2 is omitted and the embeddings in step~3 are computed with EmbeddingGemma~2 \citep{embeddinggemma2}. All features are computed with Algorithm~\ref{alg:features}; the sample split, the estimation sample, the tabular controls, the learners, and the random seeds are unchanged.}

\revised{EmbeddingGemma~2 is an open multimodal embedding model released by Google DeepMind in October 2026. Built on Gemma 4, it maps text, images, video, and audio into a shared 768-dimensional space. Of its 740M parameters, we use the 270M-parameter text model and the 170M-parameter vision encoder. The description and the image of each product are encoded jointly, with the task prefix for clustering (checkpoint \texttt{google/embeddinggemma-2}, revision \texttt{<<revision>>}). As the model is trained with Matryoshka representation learning, we keep the leading 256 coordinates of each embedding and re-normalize them instead of applying the random projection in Algorithm~\ref{alg:features}. Unlike the fine-tuned embeddings, these embeddings are not trained on our data and use neither tabular inputs nor the price and quantity signals. The comparison therefore changes the encoders, their inputs, and the fine-tuning at the same time.}

\paragraph{Clusters.} <<cluster_text>>

\begin{figure}[!htbp]
\centering
\linespread{1}\selectfont
\RevisionBlue
\footnotesize
\begin{tabular}{>{\centering\arraybackslash}m{0.12\textwidth} >{\centering\arraybackslash}m{0.52\textwidth}}
<<cluster_figure_rows>>
\end{tabular}
\caption{\revised{Products closest to the five cluster centers (EmbeddingGemma~2, Text + Image). The clusters define the similarity features $X^{sim}_i$ used in Table~\ref{tab:gemma_r2} and Figures~\ref{fig:gemma_cate} and \ref{fig:gemma_sorted}.}}
\label{fig:gemma_clusters}
\end{figure}

\paragraph{Prediction.} <<r2_text>>

\begin{table}[!htbp]
\centering
\linespread{1}\selectfont
\RevisionBlue
\small
\caption{\revised{Test $R^2$ scores (in \%) for predicting quantity and price signals: fine-tuned (FT) vs.\ EmbeddingGemma~2 (EG) embeddings.}}
\label{tab:gemma_r2}
\resizebox{1\textwidth}{!}{
\begin{tabular}{lcccccccc}
\toprule
 & \multicolumn{2}{c}{$Q_{it}$} & \multicolumn{2}{c}{$P_{it}$} & \multicolumn{2}{c}{$\Delta Q_{it}$} & \multicolumn{2}{c}{$\Delta P_{it}$}\\
\cmidrule(lr){2-3}\cmidrule(lr){4-5}\cmidrule(lr){6-7}\cmidrule(lr){8-9}
Method [features] & FT & EG & FT & EG & FT & EG & FT & EG\\
\midrule
<<table_r2>>
\bottomrule
\end{tabular}
}
\caption*{\footnotesize \revised{Note: All models are trained on a training set and scores are evaluated on a test set. Predictions use lagged values of time-varying controls. Negative $R^2$ values are set to $0$. The FT columns repeat the corresponding results of the main text.}}
\end{table}

\paragraph{Homogeneous elasticities.} <<ate_text>>

\begin{table}[!htbp]
\centering
\linespread{1}\selectfont
\RevisionBlue
\caption{\revised{Estimated inverse-rank elasticities based on the partially linear dynamic model: fine-tuned vs.\ EmbeddingGemma~2 embeddings.}}
\label{tab:gemma_ate}
\resizebox{1\textwidth}{!}{
\begin{tabular}{lcccccc}
\toprule
 & \multicolumn{3}{c}{Fine-tuned embeddings} & \multicolumn{3}{c}{EmbeddingGemma~2}\\
\cmidrule(lr){2-4}\cmidrule(lr){5-7}
Specification of Control Function (State $S_{t}$) & coef & std err & 90\% CI & coef & std err & 90\% CI\\
\midrule
<<table_ate>>
\bottomrule
\end{tabular}
}
\caption*{\footnotesize \revised{Note: Standard errors are clustered at the product level. The fine-tuned columns repeat the estimates of the main text. Model I-1 without $X^{e}$ or $X^{sim}$ does not use embeddings.}}
\end{table}

\paragraph{Heterogeneous elasticities.} <<het_text>>

\begin{figure}[!htbp]
\centering
\linespread{1}\selectfont
\RevisionBlue
\includegraphics[width=.72\textwidth]{<<figure_dir>>/cate_estimates_comparison.png}
\caption{\revised{Parameters of the estimated inverse-rank elasticity function and 90\% confidence intervals: fine-tuned (top) vs.\ EmbeddingGemma~2 (bottom) embeddings.}}
\caption*{\footnotesize \revised{Note: Covariance is clustered at the product level. The LR models are estimated using OLS, the PLR model using DML with cross-fitted boosted trees. The top panel shows the estimates of the corresponding figure of the main text. The cluster similarities refer to different clusters for the two embeddings.}}
\label{fig:gemma_cate}
\end{figure}

\begin{figure}[!htbp]
\centering
\linespread{1}\selectfont
\RevisionBlue
\begin{subfigure}{\textwidth}
\centering
\includegraphics[width=.7\textwidth]{<<figure_dir>>/sorted_effects_finetuned.png}
\caption{\revised{Fine-tuned embeddings}}
\end{subfigure}
\begin{subfigure}{\textwidth}
\centering
\includegraphics[width=.7\textwidth]{<<figure_dir>>/sorted_effects_gemma.png}
\caption{\revised{EmbeddingGemma~2}}
\end{subfigure}
\caption{\revised{Sorted inverse-rank elasticity and 90\% pointwise confidence bands: fine-tuned vs.\ EmbeddingGemma~2 embeddings.}}
\caption*{\footnotesize \revised{Note: Panel (a) repeats the corresponding figure of the main text.}}
\label{fig:gemma_sorted}
\end{figure}

\paragraph{Sensitivity analysis.} <<sensitivity_text>>

\begin{table}[!htbp]
\centering
\linespread{1}\selectfont
\RevisionBlue
\caption{\revised{Omitted-confounding scenarios at the residual shares of Table~\ref{tab:sensitivity_calibrated}: fine-tuned vs.\ EmbeddingGemma~2 embeddings.}}
\label{tab:gemma_sensitivity}
\begin{tabular}{lccccc}
\toprule
 & & \multicolumn{2}{c}{Bias bounds} & \multicolumn{2}{c}{Confidence bounds}\\
\cmidrule(lr){3-4}\cmidrule(lr){5-6}
Scenario & Residual share & Lower & Upper & Lower & Upper\\
\midrule
\multicolumn{6}{l}{\textit{Fine-tuned embeddings} ($\hat\alpha=<<alpha_main>>$)}\\
<<table_sens_main>>
\midrule
\multicolumn{6}{l}{\textit{EmbeddingGemma~2} ($\hat\alpha=<<alpha_gemma>>$)}\\
<<table_sens_gemma>>
\bottomrule
\end{tabular}
\caption*{\footnotesize \revised{Note: $R_Y^2=R_D^2$ equal to the residual share, with worst-case alignment $|\rho|=1$. The residual shares are the calibration of the main analysis ($r_0=<<r0_main_7>>$) and its 1.5 and 2 multiples, for both embeddings. The bounds refer to the inverse-rank elasticity of Model I-3 with embeddings. Bias bounds allow for omitted confounding; confidence bounds additionally account for sampling uncertainty (one-sided 90\% bounds). The fine-tuned rows repeat Table~\ref{tab:sensitivity_calibrated}.}}
\end{table}

<<conclusion_text>>
\par\endgroup
